% ===================================================================
% ABSTRACT (ENGLISH)
% ===================================================================
\chapter*{Abstract}
\addcontentsline{toc}{chapter}{Abstract}

\textbf{Background.} Federated learning allows intrusion detection systems for Internet-of-Things
networks to be trained without moving raw traffic off the devices that observe it, but it
raises a maintenance problem: when a node is identified as compromised, or when a right to
erasure is exercised, the influence of its data must be removed from the model. Deleting the
raw records is not sufficient, because the model retains what it learned from them in its
weights. The federated unlearning literature expresses the cost of removal in epochs and
communication rounds, and measures it in simulation. On a real edge device --- constrained
simultaneously by compute, heat, storage and power delivery --- that abstraction is not
enough. This thesis measures the \textbf{physical cost of unlearning} on real hardware.

\textbf{Methods.} A hybrid four-client federation was built on Flower 1.29 in deployment mode,
with the aggregator and three simulated nodes running on a host machine and the fourth node
being a physical Raspberry Pi 5 (4 GB) joining over the local network. The device is
deliberately constrained: its fan is unplugged, which forces thermal throttling, and its
storage is a class A1 SD card. Power is measured with an external FNIRSI FNB58 meter placed
inline on the supply. Using the TON\_IoT dataset with 1,000,000 rows per node, two approaches
to exact local unlearning are compared: \textbf{naive retraining} from scratch on the retained
data, and \textbf{SISA unlearning} ($S = 5$ shards × $R = 5$ slices) with rollback and replay of
only the affected slices. Removal is triggered by a label-flipping poison confined to slices
3–4 of a single shard, which defines a well-specified set of roughly 63,000 samples that must
be forgotten. The design was pre-registered and frozen before measurement, paired by seed,
with $N = 10$ pairs (seeds 42–51) and counterbalanced arm order. Statistics are
non-parametric: Wilcoxon signed-rank tests, Cliff's $\delta$ and bootstrap 95\,\% confidence
intervals.

\textbf{Results.} SISA reduces time-to-recovery from a median of 242.1 s to 29.5 s (\textbf{8.2×}), net
recovery energy from 0.093 Wh to 0.017 Wh (\textbf{5.4×}), and time spent thermally throttled from
233 s to 27.5 s (\textbf{8.5×}). All three reach $p = 0.0020$ --- the two-sided floor at $N = 10$ --- with Cliff's $\delta = +1.00$, that is, complete separation: no SISA run overlaps any naive
run. Cost does shift toward storage (12.05 MB written versus 2.75 MB), but the absolute
magnitude is negligible for a 12,865-parameter model. SISA's training-phase overhead is
43.67 MB and a median of 5.6 s of checkpoint I/O, which a single recovery repays more than
thirtyfold. \textbf{Utility requires careful reading.} Re-scoring every saved global model on a
held-out 100,000-sample test set that is 96.5\,\% attack shows that at the deployed 0.5
threshold the SISA-recovered model predicts every flow as an attack: specificity 0.000,
balanced accuracy 0.500 (chance) and MCC 0.017, against 0.237, 0.598 and 0.177 for naive
retraining; its recall of 1.000 and $F_1$ of 0.982 are trivial artefacts of class imbalance.
The collapse is, however, \textbf{a matter of calibration}: with the threshold tuned, SISA reaches
a balanced accuracy of \textbf{0.709 against 0.671} for naive retraining, better on 9 of 10 seeds
and better on specificity and MCC alike. It precedes recovery and localises to the parameter
averaging of the constituent models --- individually they discriminate, the average of their
weights does not.

\textbf{Conclusion.} SISA makes unlearning on an edge device eightfold cheaper, at a negligible
standing cost and \textbf{without sacrificing utility --- provided the decision threshold is
calibrated}. Left uncalibrated it delivers a detector that labels everything an attack: a
failure the conventional recall and $F_1$ metrics conceal entirely under severe class
imbalance, and one that originates in the aggregation rather than in the unlearning. The
contribution of this thesis is a methodology for measuring the physical cost of unlearning,
and an empirical basis for choosing between the two approaches, expressed in the quantities a
systems designer actually pays.

\textbf{Keywords:} federated learning, machine unlearning, SISA, intrusion detection, Internet of
Things, edge computing, energy consumption, thermal throttling, class imbalance, right to be
forgotten.

